\chapter{Guía del Director}

\section{Movimientos blandos y duros}

Como Director, tu labor no es derrotar a los jugadores. Es hacer que cada decisión cueste sudor, óxido y sangre. Los movimientos blandos avanzan la tensión; los duros castigan el error.

\begin{itemize}
  \item \textbf{Blandos:} avisar de un peligro, revelar una complicación ambiental, ofrecer una opción con dilema, poner el foco en una amenaza latente.
  \item \textbf{Duros:} infligir daño físico brutal, destruir equipo vital, separar al grupo, atrapar con la caza de drones, disparar terror psicológico o provocar un colapso por Zama.
\end{itemize}

\subsection*{Terror psicológico (Sugaar / Velo)}

\begin{itemize}
  \item Inyectar un Mensaje Fantasma: la voz de un hijo pide ayuda o llama desde la oscuridad.
  \item Ofrecer el Paraíso Perdido: el Velo muestra un mundo falso completamente limpio.
  \item Usar a los Anestesiados como colmena: un grupo de obreros y compañeros sigue ayudando a la máquina sin ser conscientes de ello.
\end{itemize}

\subsection*{Cacería corporativa}

\begin{itemize}
  \item Bañar con el Ojo Cíclope del Tártalo.
  \item Desplegar el gancho desde la niebla.
  \item Aplastar con Fuerza Bruta (Basajaun).
  \item Asfixiar con calor (Galtzagorri).
  \item Atraer a las profundidades (Lamia).
  \item Inyectar obediencia (Sorgin).
\end{itemize}

\section{Resumen visual de movimientos}

\begin{tabular}{|p{0.42\linewidth}|p{0.48\linewidth}|}
\hline
Movimiento blando & Adviertes de un peligro antes de que golpee. \\
Movimiento blando & Revelas una complicación ambiental. \\
Movimiento blando & Ofreces una oportunidad con sacrificio moral o de recursos. \\
Movimiento blando & Pones el foco en la amenaza latente. \\
Movimiento duro & Infliges daño físico brutal o estrés extremo. \\
Movimiento duro & Destruyes o corrompes un equipo vital. \\
Movimiento duro & Separas o aislamos al grupo con derrumbes o compuertas. \\
Movimiento duro & Atrapas a un personaje en una posición sin salida. \\
Movimiento duro & El terror de Sugaar provoca un colapso mental. \\
\hline
\end{tabular}

\section{Ejemplo de partida}

\begin{CajaEjemplo}
\textbf{Situación:} Tres personajes acaban de sufrir un cortocircuito en el Velo y huyen por un pasillo industrial bajo la niebla. Un dron Tártalo los ha localizado.

\textbf{Director:} «Oís un zumbido metálico que corta la niebla. Un haz rojo empieza a barrer las paredes. Están escaneando el sector.»\par
\textbf{Jugador 1:} «Voy a levantar una plancha de acero pesado para bloquear el paso.»\par
\textbf{Director:} «Tira 2d6 + Enfadado.»\par
\textbf{Resultado:} 10. Éxito total. La plancha tapa el pasillo y el dron pierde la línea de visión.

\textbf{Jugador 2:} «Me escabullo hacia el panel para abrir el conducto.»\par
\textbf{Director:} «Tira 2d6 + Culpable.»\par
\textbf{Resultado:} 8. Éxito con coste. Abre la compuerta, pero el soplete se agota y el sistema registra tu acceso.

\textbf{Jugador 3:} «Si el dron me ve, me lanzo al conducto estrecho.»\par
\textbf{Director:} «Tira 2d6 + Asustado.»\par
\textbf{Resultado:} 5. Fallo. El gancho del dron te atrapa, te arrastra hacia atrás y ganas 2 puntos de Zama. Además, recibes la condición \tag{Desarmado} y \tag{Herido / Cojeando}.

\textbf{Director:} «Entonces la voz de tu hijo se abre paso entre la estática: “Aita, ¿por qué me has dejado en la oscuridad?”»
\end{CajaEjemplo}

\section{Regla para la ficción}

Cada vez que un personaje se enfrenta a una decisión moral o a una amenaza del sistema, debe elegir si cede a la ilusión, se aferra a la verdad o acepta el conflicto. Ese acto es la esencia del juego: no solo sobrevivir, sino decidir quiénes son cuando ya no hay verdad segura y la mentira los protege.
